\documentclass[]{article}

\usepackage{amsmath}
\usepackage{multirow}
\usepackage{natbib}
\usepackage{graphicx} 


\begin{document}

The damped harmonic oscillator is a cornerstone model in science and engineering, describing phenomena ranging from the vibrations of civil structures to the dynamics of micro-electromechanical systems. The predictive power of these models is critical for applications such as structural health monitoring, sensor design, and system control. However, the fidelity of any such model is fundamentally limited by the precision of its physical parameters, namely mass, stiffness, and damping. In practice, these parameters are determined through measurement and are thus subject to uncertainty, which can propagate through the model and lead to significant discrepancies between predicted and observed system behavior.

A subtle but critical challenge in modeling these systems is that the influence of parameter uncertainty is not constant over time. As an oscillator evolves and dissipates energy, the sensitivity of its state to different parameters changes. For instance, an error in the system's mass may dominate the energy budget during the initial, high-amplitude oscillations, whereas an error in the damping coefficient may have a more cumulative effect that becomes prominent only at later times. This temporal evolution of parameter influence creates a practical dilemma for system identification: it is unclear which time windows are optimal for reliably estimating a specific parameter. Without a quantitative understanding of these dynamics, experimental design and model validation efforts may be inefficient or misleading.

This paper addresses this challenge by quantifying the temporal limits of parameter identifiability in underdamped harmonic oscillators. We investigate how the sensitivity of the system's mechanical energy, described by the equation $E(t) = \frac{1}{2} m v(t)^2 + \frac{1}{2} k x(t)^2$, responds to perturbations in mass ($m$) and the damping coefficient ($b$). By employing a Jacobian-based sensitivity analysis, we compute the partial derivatives of the energy with respect to each parameter at every time step. This approach allows us to track the time-varying contribution of each parameter to the uncertainty in the system's energy, thereby mapping the evolving landscape of parameter influence directly from the system's dynamics.

Our analysis reveals a distinct and rapid transition in which the dominant source of energy uncertainty shifts from mass to the damping coefficient. We define this critical crossover point as the "Information Horizon," a temporal boundary beyond which the system's dynamics contain more information about dissipative processes than inertial properties. Through the study of a population of simulated oscillators, we demonstrate a systematic relationship between the timing of this horizon and the system's damping ratio. This work establishes a quantitative framework for understanding the time-dependent nature of parameter sensitivity, providing clear insights into the operational limits for robust parameter estimation and model prediction in dissipative systems.

\end{document}
                